\newcommand{\memoryChart}[2][]{%
    % #1 = optional max difficulty (empty -> show all)
    % #2 = CSV file path

    % Load the CSV data
    \pgfplotstableread[col sep=comma]{#2}\loadedMemoryData

    % Store the max difficulty in a new macro
    \ifx\relax#1\relax
        \def\maxDifficulty{0}% Use 0 or some other indicator for no max
        \def\hasMax{0}%
    \else
        \edef\maxDifficulty{#1}% <--- *** Crucial change: Store #1's content ***
        \def\hasMax{1}%
    \fi

    \ifnum\hasMax=1
        % If max is set, \currentXmax is the number
        \edef\currentXmax{\maxDifficulty}
    \else
        % If no max is set, \currentXmax is the pgfplots 'no value' token
        \let\currentXmax\pgfkeysnovalue
    \fi

    \begin{tikzpicture}

        % --- Left axis (ILP memory) ---
        \begin{axis}[
                title={Hauptspeicher (Bytes)},
                name=ilpaxis,
                axis y line*=left,
                xlabel={Difficulty},
                ylabel={ILP Speichernutzung (Bytes)},
                ymode=log,
                % xmax=\currentXmax,
                ymajorgrids=true,
                width=0.95\textwidth,
                height=7cm,
                grid=major,
                legend style={at={(0.5,-0.2)}, anchor=north, legend columns=-1},
                table/x=difficulty,
            ]

            % --- ILP Memory Usage ---
            \addplot[
                mark=*,
                color=blue,
                line width=1pt,
                smooth,
                error bars/.cd,
                y dir=both,
                y explicit,
                /pgfplots/x filter/.code={\ifnum\hasMax=1 \ifdim\pgfmathresult pt>\maxDifficulty pt \def\pgfmathresult{} \fi \fi}
            ] table [
                    x=difficulty,
                    y=ILP,
                    y error expr=\thisrow{ILP stdev},
                ] {\loadedMemoryData};
            \addlegendentry{ILP}

            % --- ILP GPU RAM ---
            \addplot[
                mark=*,
                color=purple,
                line width=1pt,
                smooth,
                error bars/.cd,
                y dir=both,
                y explicit,
                /pgfplots/x filter/.code={\ifnum\hasMax=1 \ifdim\pgfmathresult pt>\maxDifficulty pt \def\pgfmathresult{} \fi \fi}
            ] table [
                    x=difficulty,
                    y={ILP GPU},
                    y error={ILP GPU stdev},
                ] {\loadedMemoryData};
            \addlegendentry{ILP GPU RAM}

            % --- Hill-Climbing v1 ---
            \addplot[
                mark=square*,
                color=red,
                line width=1pt,
                smooth,
                error bars/.cd,
                y dir=both,
                y explicit,
                /pgfplots/x filter/.code={\ifnum\hasMax=1 \ifdim\pgfmathresult pt>\maxDifficulty pt \def\pgfmathresult{} \fi \fi}
            ] table [
                    x=difficulty,
                    y={Hill Climbing v1},
                    y error expr=\thisrow{Hill Climbing v1 stdev},
                ] {\loadedMemoryData};
            \addlegendentry{Hill-Climbing v1}

            % --- Hill-Climbing v3 ---
            \addplot[
                mark=triangle*,
                color=green!60!black,
                line width=1pt,
                smooth,
                error bars/.cd,
                y dir=both,
                y explicit,
                /pgfplots/x filter/.code={\ifnum\hasMax=1 \ifdim\pgfmathresult pt>\maxDifficulty pt \def\pgfmathresult{} \fi \fi}
            ] table [
                    x=difficulty,
                    y={Hill Climbing v3},
                    y error expr=\thisrow{Hill Climbing v3 stdev},
                ] {\loadedMemoryData};
            \addlegendentry{Hill-Climbing v3}

            % --- Offering-Order ---
            \addplot[
                mark=diamond*,
                color=orange,
                line width=1pt,
                dashed,
                smooth,
                error bars/.cd,
                y dir=both,
                y explicit,
                /pgfplots/x filter/.code={\ifnum\hasMax=1 \ifdim\pgfmathresult pt>\maxDifficulty pt \def\pgfmathresult{} \fi \fi}
            ] table [
                    x=difficulty,
                    y={Offering Order},
                    y error expr=\thisrow{Offering Order stdev},
                ] {\loadedMemoryData};
            \addlegendentry{Offering-Order}

        \end{axis}
    \end{tikzpicture}
}

\newcommand{\vramChart}[2][]{%
    % #1 = optional max difficulty (empty -> show all)
    % #2 = CSV file path

    \pgfplotstableread[col sep=comma]{#2}\loadedVramData

    \ifx\relax#1\relax
        \def\vramMaxDifficulty{0}%
        \def\vramHasMax{0}%
    \else
        \edef\vramMaxDifficulty{#1}%
        \def\vramHasMax{1}%
    \fi

    \begin{tikzpicture}
        \begin{axis}[
                title={ILP GPU VRAM (Bytes)},
                axis y line*=left,
                xlabel={Difficulty},
                ylabel={VRAM Nutzung (Bytes)},
                ymajorgrids=true,
                width=0.95\textwidth,
                height=7cm,
                grid=major,
                legend style={at={(0.5,-0.2)}, anchor=north, legend columns=-1},
                table/x=difficulty,
            ]

            \addplot[
                mark=pentagon*,
                color=cyan!70!black,
                line width=1pt,
                dashed,
                smooth,
                error bars/.cd,
                y dir=both,
                y explicit,
                /pgfplots/x filter/.code={\ifnum\vramHasMax=1 \ifdim\pgfmathresult pt>\vramMaxDifficulty pt \def\pgfmathresult{} \fi \fi}
            ] table [
                    x=difficulty,
                    y={ILP GPU VRAM},
                    y error={ILP GPU VRAM stdev},
                ] {\loadedVramData};
            \addlegendentry{ILP GPU VRAM}

        \end{axis}
    \end{tikzpicture}
}
